\section{模型学会了怎样的“耳朵”：Inspecting the Learned Front End}

只看 mAP 仍不知道模型如何表示声音。论文进一步检查第一层 filters：对每个 learned filter 做 Fourier transform，按 center frequency 排序，再观察频带密度和 filter shape。这个过程把神经网络参数转换成信号处理读者熟悉的频率响应，是本讲最接近 mechanistic inspection 的部分。

\subsection{从 Filter Impulse Response 到 Frequency Response}

若第 $j$ 个前端 filter 的时域权重为 $h_j[n]$，其频率响应为
\[
H_j(\omega)=\sum_n h_j[n]e^{-i\omega n}.
\]
可用峰值位置定义近似 center frequency
\[
\omega_j^*=\operatorname*{argmax}_{\omega}|H_j(\omega)|.
\]
将 filters 按 $\omega_j^*$ 排序后，固定 Fourier bank 会近似沿直线均匀展开；learned bank 若在某些频段分配更多 filters，则曲线会弯曲或出现 step pattern，表示模型按任务重新分配频率 resolution。

\lecturefigure{slide-45-ideal-time-frequency.jpg}{理想固定 Fourier bank 与 learned task-specific frequency organization 的比较。}{Stanford Online 官方视频 00:44:36--00:46:09。}

\begin{knowledgebox}{读图：横轴、纵轴与曲线分别是什么}
横轴是排序后的 filter index，纵轴是每个 filter 响应最大的 center frequency。直线表示均匀频率分布；弯曲表示某些频段获得更密集的 filters。不同任务出现不同曲线，支持“front end 随任务自适应”，但曲线本身不告诉我们哪种表示在分布外更稳健。
\end{knowledgebox}

\teachervoice{Verma 把这种现象描述为计算机按任务“调整自己的耳朵”。这句比“模型学到了 spectrogram”更准确：它学的是一组非均匀、任务相关的 analysis filters，而不是复制某个固定 STFT。}

\subsection{DSP-like Filters：Onset、Window 与 Sinusoid}

将 individual filters 可视化后，可以看到类似 onset detector、windowing function 与 sinusoidal basis 的形状；同一频段还可能出现多个相位不同的 filters。它们说明 gradient descent 重新发现了若干经典 DSP building blocks，同时也组合出固定 basis 未预设的形状。

\lecturefigure{slide-46-learned-filter-bank.jpg}{Learned filters：onset-sensitive、window-like、sinusoidal 与 phase-diverse patterns。}{Stanford Online 官方视频 00:46:10--00:47:18。}

\begin{warningbox}{Visual resemblance 不是完整机制证明}
一个 filter 看起来像 Hamming window 或 sinusoid，只说明 shape/response 相似。要证明 causal function，还需 intervention、ablation、输入合成测试与跨 checkpoint 稳定性；多个 filters 也可能共同实现 distributed feature。
\end{warningbox}

\begin{importantbox}{传统 DSP 与 End-to-End Learning 的关系}
本讲最有价值的结论不是“神经网络取代信号处理”，而是二者可以形成闭环：DSP 提供可解释 basis、multi-scale transform 与诊断工具；end-to-end learning 根据任务重新分配频带和组合 filters；分析结果再反过来指导 architecture。
\end{importantbox}

\subsection{本章小结}

Fourier inspection 将 learned weights 映射到 center frequency 与 filter shape，显示前端会形成任务相关 filterbank，并重现部分经典 DSP structure。它是有力的表示证据，但距离完整 causal circuit 仍有一步。

